% ===== EXPERIMENTS =====
% ===== EXECUTION PLAN (for efficient experiment scheduling) =====
% Main-paper figures/tables to prioritize (in order):
%   P0  Fig-ZS-Curve (must-have):
%       Metric vs c where c scales parallel component, c in {-2,-1,0,0.5,1,2,4}.
%       Three settings: Attn-only / MLP-only / All.
%       Purpose: establish "inference null direction" (not only deletion-safe).
%       Suggested artifact:
%         results/paper/zeroshot_scale_sweep_{model}.json
%         figures/paper/fig_zeroshot_scale_curve.pdf
%
%   P1  Fig-Profile:
%       Layerwise PR^l depth profile for Attn/MLP/Combined over models.
%       Suggested artifact:
%         results/paper/profile_layerwise_summary.json
%         figures/paper/fig_profile_layerwise.pdf
%
%   P2  Fig-Pretrain-Loss:
%       Training loss curves: baseline / residual-space control / value-space control.
%       Suggested artifact:
%         results/paper/pretrain_loss_curves.csv
%         figures/paper/fig_pretrain_loss_main.pdf
%
%   P3  Tab-Pretrain-Summary:
%       Final ppl + diagonal concentration + off-diagonal entropy.
%       Suggested artifact:
%         results/paper/pretrain_final_metrics.csv
%
%   P4  Ablation (can be one compact figure):
%       Gate-strength sweep and/or block-target sweep.
%       Suggested artifact:
%         results/paper/ablation_gate_sweep.csv
%         figures/paper/fig_ablation_gate_sweep.pdf
%
% Move to appendix (do not block main writing):
%   A1  Per-model breakdown tables for all zero-shot tasks.
%   A2  Extra c-grid points and seed-by-seed curves.
%   A3  Calibration size sensitivity and scale transfer.
%   A4  Extra attention-map galleries.
%
% Narrative flow:
%   §4.1 Setup: models, calibration corpus, benchmarks for both zero-shot and pretraining.
%   §4.2 Profiles: what R^l looks like across models/blocks/depths.
%         Key: both Attn and MLP are affected; ratio grows with depth; Attn > MLP.
%   §4.3 Zero-shot removal: parallel component is redundant at inference.
%         Key: small perf drop (<X%) for all three variants.
%   §4.4 Pretraining: removing parallel component during training helps.
%         Key: compare residual-space and value-space controls.
%   §4.5 [method results placeholder].
%   §4.6 Ablation: which blocks? calibration size? scale?

\section{Experiments}
\label{sec:experiments}

\subsection{Experimental Setup}
\label{sec:setup}

\paragraph{Models for zero-shot analysis.}
We measure $\PR{l}$ on a diverse set of publicly available pretrained transformers spanning
decoder-only and encoder-only architectures: [e.g., Llama-3-8B,
Llama-3-70B, Mistral-7B-v0.3, Qwen2-7B, Qwen2-72B, and encoder-only models
such as RoBERTa-large and DeBERTa-v3-large].
This covers a range of model sizes ([RANGE]M--[RANGE]B parameters),
training corpora, and positional encoding schemes (RoPE, ALiBi, learned
absolute), enabling us to assess the generality of the parallel component.

\paragraph{Calibration corpus.}
For zero-shot $\PR{l}$ measurement and parallel control, we use
a held-out calibration corpus sampled uniformly from [CORPUS] (e.g., a held-out slice of
C4 or The Pile), tokenized to a maximum of 512 tokens.
We verify that results are stable across corpus choices in \S\ref{sec:ablation}.

\paragraph{Zero-shot evaluation benchmarks.}
We evaluate the effect of zero-shot parallel component removal on:
(a) causal language modeling perplexity on [HELD-OUT CORPUS] for decoder-only
models; and (b) [DOWNSTREAM TASKS — e.g., MMLU, HellaSwag, WinoGrande, ARC]
using zero-shot accuracy.
All evaluations use the \texttt{lm-evaluation-harness}~\citep{gao2021framework}.

\paragraph{Pretraining setup.}
For pretraining experiments, we follow a two-stage scale protocol.
\textbf{Stage 1 (small-scale validation):} we start with GPT-style models
below 500M parameters to rapidly test intervention behavior and tune control
strength.
\textbf{Stage 2 (scale extension):} we extend the same intervention settings
to larger models up to 2.7B parameters to test whether the observed trends
persist with scale.
In both stages we compare: (1)~baseline, (2)~residual-space parallel control,
(3)~value-space parallel control, and (4)~gated parallel control variants.
Models in each stage share architecture family, optimizer, and learning-rate
schedule, so differences are attributable to control strategy rather than
training recipe changes. Hyperparameters are in
Appendix~\ref{app:pretraining}.

\begin{table*}[t]
  \caption{Model configurations, hyperparameters, and validation loss across
  parallel-control methods. We fix architecture and optimization settings within each scale;
  validation-loss entries are filled after runs complete.}
  \label{tab:train-config-val-loss}
  \centering
  \small
  \begin{tabular}{cccccccc|cccc}
    \toprule
    \#Act.\ Params & Tokens & $L_b$ & $H$ & $d_{\text{model}}$ & $d_{\text{ff}}$ & lr & batch size &
    \multicolumn{4}{c}{Val.\ Loss} \\
    \cmidrule(lr){9-12}
    & & & & & & & &
    Baseline & Residual-space control & Value-space control & Gated control \\
    \midrule
    194M & TBD & 12 & 14 & 896 & 3584 & $2\times 10^{-4}$ & 32 &
    TBD & TBD & TBD & TBD \\
    241M & TBD & 14 & 15 & 960 & 3840 & $2\times 10^{-4}$ & 32 &
    TBD & TBD & TBD & TBD \\
    296M & TBD & 14 & 16 & 1024 & 4096 & $2\times 10^{-4}$ & 32 &
    TBD & TBD & TBD & TBD \\
    436M & TBD & 16 & 18 & 1152 & 4608 & $2\times 10^{-4}$ & 16 &
    TBD & TBD & TBD & TBD \\
    528M & TBD & 18 & 20 & 1280 & 5120 & $2\times 10^{-4}$ & 16 &
    TBD & TBD & TBD & TBD \\
    \bottomrule
  \end{tabular}
\end{table*}

% ------------------------------------------------------------------
\subsection{Parallel Component Profiles}
\label{sec:profiles}

Figure~\ref{fig:profiles} shows $\PR{l}$ for Attn-only, MLP-only, and
combined residual updates, averaged across all evaluated models.

Three consistent patterns emerge across all models.
\textbf{First, the parallel component is substantial and ubiquitous:}
even at the earliest layers, $\PR{l}$ exceeds [X]\%, meaning a significant
fraction of each block's residual update is directionally redundant from the
start of the network.
\textbf{Second, $\PR{l}$ grows monotonically with depth}
for both Attn and MLP blocks, consistent with the structural argument in
\S\ref{sec:why}: as the accumulated residual norm grows, the same
bounded block output becomes progressively more aligned with the large
pre-existing vector.
\textbf{Third, Attn blocks consistently exhibit higher $\PR{l}$ than MLP
blocks at the same depth.}
We explain this asymmetry mechanistically in \S\ref{sec:mechanism}:
the attention diagonal $\Adiag{t}$, which is structurally driven toward
high values by the residual norm growth, produces a self-referential
contribution that MLP blocks lack.
Appendix~\ref{app:profiles} shows per-model profiles.

\paragraph{Transition to zero-shot intervention.}
Section~\ref{sec:profiles} establishes where parallel components are strong.
We next test their functional role directly in zero-shot settings by
intervening on the parallel part itself: first by hard deletion, then by
deletion-ratio ablation, and finally by robustness analysis.

% ------------------------------------------------------------------
\subsection{Zero-Shot Parallel Control}
\label{sec:zeroshot}

\paragraph{Step 1: Hard deletion.}
We first apply hard deletion of the parallel component at inference time:
\begin{equation}
  \tilde{\Du{l}{t}} = \Duperp{l}{t}.
\end{equation}
We run this intervention under three variants:
\textbf{(A)} Attn-only, \textbf{(B)} MLP-only, and \textbf{(C)} all sub-layers.
Table~\ref{tab:zeroshot} summarizes the resulting task deltas.

\begin{figure}[t]
  \centering
  % Replace with:
  % \includegraphics[width=0.96\linewidth]{figures/paper/fig_zeroshot_parallel_perp_sweep.pdf}
  \fbox{\rule{0pt}{5.0cm}\rule{0.92\linewidth}{0pt}}
  \caption{\textbf{Zero-shot ratio ablation (main result).}
           Performance under varying deletion ratio of the parallel component
           (left) and, for comparison, the perpendicular component (right).
           Curves correspond to Attn-only, MLP-only, and all-sub-layer
           interventions. Parallel deletion remains broadly robust over a wide
           ratio range, while perpendicular deletion is substantially more
           sensitive.}
  \label{fig:zeroshot-scale-sweep}
\end{figure}

\begin{table}[t]
  \caption{Zero-shot performance before and after hard parallel deletion.
           ``PPL'' is perplexity (lower is better);
           downstream scores are zero-shot accuracy (\%, higher is better).
           All changes are relative to the unmodified baseline.
           \textdagger: deletion applied to Attn sub-layers only (variant A);
           \textdaggerdbl: MLP only (B); $\S$: all sub-layers (C).}
  \label{tab:zeroshot}
  \centering
  \small
  \begin{tabular}{lrrrrr}
    \toprule
    Model & PPL $\Delta$ & HellaSwag $\Delta$ & MMLU $\Delta$ & WinoGrande $\Delta$ & ARC $\Delta$ \\
    \midrule
    \multicolumn{6}{l}{\textit{Attn-only removal (variant A)}} \\
    {[MODEL A]} & $+$-- & $-$-- & $-$-- & $-$-- & $-$-- \\
    {[MODEL B]} & $+$-- & $-$-- & $-$-- & $-$-- & $-$-- \\
    \midrule
    \multicolumn{6}{l}{\textit{MLP-only removal (variant B)}} \\
    {[MODEL A]} & $+$-- & $-$-- & $-$-- & $-$-- & $-$-- \\
    {[MODEL B]} & $+$-- & $-$-- & $-$-- & $-$-- & $-$-- \\
    \midrule
    \multicolumn{6}{l}{\textit{All sub-layers (variant C)}} \\
    {[MODEL A]} & $+$-- & $-$-- & $-$-- & $-$-- & $-$-- \\
    {[MODEL B]} & $+$-- & $-$-- & $-$-- & $-$-- & $-$-- \\
    \bottomrule
  \end{tabular}
\end{table}

\paragraph{Step 2: Deletion-ratio ablation.}
Beyond hard deletion, we sweep a ratio parameter $\gamma\in[0,1]$ and define
\begin{equation}
  \tilde{\Du{l}{t}} = (1-\gamma)\Dup{l}{t} + \Duperp{l}{t},
\end{equation}
where $\gamma=0$ recovers baseline and $\gamma=1$ is full parallel deletion.
Figure~\ref{fig:zeroshot-scale-sweep} shows that performance remains largely
flat across a broad range of $\gamma$, reinforcing that inference is robust to
parallel suppression.

\paragraph{Step 3: Why is this robust?}
We analyze robustness from an architectural perspective. Parallel updates are
direction-preserving and are strongly attenuated by subsequent normalization,
while directional changes induced by $\Duperp{l}{t}$ alter token-to-token
relative order and therefore more directly affect prediction boundaries and
generation ranking. This explains why perpendicular interventions are far more
sensitive than parallel interventions in zero-shot evaluation.

\begin{table}[t]
  \caption{General-task zero-shot performance with per-task breakdown:
  baseline, attention removal, and value-space parallel control.
  Scores are from \texttt{general.csv}, rounded to one decimal place.}
  \label{tab:general-zeroshot}
  \centering
  \scriptsize
  \setlength{\tabcolsep}{4pt}
  \begin{tabular}{lrrrrrr}
    \toprule
    \multirow{2}{*}{Task} & \multicolumn{3}{c}{Qwen3-1.7B} & \multicolumn{3}{c}{Qwen3-30B-A3B} \\
    \cmidrule(lr){2-4}\cmidrule(lr){5-7}
      & Baseline & Attn Rem. & XSA & Baseline & Attn Rem. & XSA \\
    \midrule
    OpenBookQA & 37.2 & 37.2 & 36.4 & 45.4 & 43.0 & 45.2 \\
    PIQA & 72.1 & 72.0 & 72.1 & 80.2 & 78.8 & 80.8 \\
    RTE & 70.0 & 66.1 & 70.8 & 82.0 & 79.1 & 82.7 \\
    Winogrande & 60.7 & 55.8 & 59.8 & 72.3 & 66.6 & 72.2 \\
    BoolQ & 77.8 & 71.2 & 77.6 & 88.7 & 88.0 & 88.6 \\
    ARC-Challenge & 53.2 & 50.2 & 53.1 & 70.0 & 64.9 & 69.9 \\
    HellaSwag & 60.4 & 58.9 & 60.4 & 77.8 & 73.6 & 77.8 \\
    MMLU & 35.0 & 23.0 & 35.0 & 67.0 & 50.0 & 67.0 \\
    GSM8K & 68.0 & 0.2 & 67.7 & 89.8 & 39.9 & 89.8 \\
    HumanEval & 38.4 & 25.0 & 42.1 & 32.3 & 14.6 & 31.7 \\
    NQ-Open & 9.7 & 7.2 & 9.7 & 25.1 & 16.4 & 25.2 \\
    DROP & 7.3 & 5.2 & 7.3 & 5.5 & -- & 5.4 \\
    \midrule
    Avg & 50.9 & 37.6 & 51.1 & 64.5 & 55.9 & 62.2 \\
    \bottomrule
  \end{tabular}
\end{table}

\begin{table*}[t]
  \caption{nanoGPT downstream evaluation results of XSA vs.\ baseline.
  We report the currently stable 0.7B results and leave 1.4B/2.7B blank until
  those runs are cleaned and revalidated. The current nanoGPT evaluation set
  uses RTE in place of SocIQA.}
  \label{tab:nanogpt-xsa}
  \centering
  \scriptsize
  \setlength{\tabcolsep}{4pt}
  \begin{tabular}{llrrrrrrrrrr}
    \toprule
    Model & Setting & ARC-E & BoolQ & HSwag & LAMBADA & OBQA & PIQA & RTE & WinoGr & Avg & $\Delta$Avg \\
    \midrule
    \multirow{2}{*}{0.7B}
      & Baseline & 39.31 & 58.35 & 31.36 & 22.30 & 30.20 & 60.72 & 51.99 & 51.30 & 46.18 & -- \\
      & XSA & 31.52 & 48.75 & 31.82 & 36.08 & 30.20 & 57.56 & 51.62 & 53.43 & 43.56 & $-2.62$ \\
    \midrule
    \multirow{2}{*}{1.4B}
      & Baseline & -- & -- & -- & -- & -- & -- & -- & -- & -- & -- \\
      & XSA & -- & -- & -- & -- & -- & -- & -- & -- & -- & -- \\
    \midrule
    \multirow{2}{*}{2.7B}
      & Baseline & -- & -- & -- & -- & -- & -- & -- & -- & -- & -- \\
      & XSA & -- & -- & -- & -- & -- & -- & -- & -- & -- & -- \\
    \bottomrule
  \end{tabular}
\end{table*}

\begin{table}[t]
  \caption{Long-context RULER performance with task-family and context-length breakdown
  (Llama-3.2-3B), comparing baseline, branch-specific removal, and XSA.
  Scores are from \texttt{ruler.csv}, rounded to one decimal place.}
  \label{tab:longcontext-ruler}
  \centering
  \scriptsize
  \setlength{\tabcolsep}{3.5pt}
  \begin{tabular}{llrrrrrr}
    \toprule
    Family & Setting & 4k & 8k & 16k & 32k & 64k & 128k \\
    \midrule
    \multirow{5}{*}{NIAH}
      & Baseline & 99.3 & 97.6 & 93.9 & 90.3 & 85.7 & 83.9 \\
      & Attn Rem. & 96.5 & 95.3 & 90.4 & 85.8 & 81.2 & 77.5 \\
      & MLP Rem. & 72.4 & 61.1 & 56.8 & 52.0 & 49.6 & 43.5 \\
      & Both Rem. & 84.8 & 71.5 & 62.6 & 55.5 & 53.0 & 48.0 \\
      & XSA & 99.4 & 97.7 & 94.0 & 90.7 & 85.9 & 83.9 \\
    \midrule
    \multirow{5}{*}{QA}
      & Baseline & 55.0 & 49.6 & 49.0 & 48.1 & 41.8 & 39.9 \\
      & Attn Rem. & 45.7 & 40.5 & 42.3 & 42.8 & 36.9 & 35.8 \\
      & MLP Rem. & 40.4 & 32.0 & 28.0 & 26.5 & 20.4 & 16.3 \\
      & Both Rem. & 44.1 & 36.7 & 31.3 & 31.6 & 24.7 & 20.3 \\
      & XSA & 54.6 & 49.4 & 49.3 & 48.5 & 42.2 & 39.6 \\
    \midrule
    \multirow{5}{*}{VT}
      & Baseline & 96.8 & 93.6 & 91.6 & 91.2 & 87.3 & 82.9 \\
      & Attn Rem. & 71.0 & 70.2 & 69.3 & 69.6 & 71.2 & 63.5 \\
      & MLP Rem. & 39.6 & 44.0 & 44.8 & 22.2 & 16.2 & 6.1 \\
      & Both Rem. & 57.7 & 57.8 & 61.4 & 40.3 & 28.8 & 9.8 \\
      & XSA & 96.1 & 93.0 & 91.4 & 91.0 & 87.8 & 83.2 \\
    \midrule
    \multirow{5}{*}{CWE}
      & Baseline & 48.0 & 11.3 & 5.8 & 1.1 & 0.6 & 0.9 \\
      & Attn Rem. & 18.6 & 8.3 & 5.3 & 2.0 & 0.5 & 0.3 \\
      & MLP Rem. & 20.6 & 8.6 & 4.8 & 2.9 & 1.9 & 2.5 \\
      & Both Rem. & 49.1 & 30.6 & 15.0 & 6.0 & 3.5 & 2.8 \\
      & XSA & 44.0 & 9.8 & 5.4 & 1.0 & 0.6 & 0.9 \\
    \midrule
    \multirow{5}{*}{FWE}
      & Baseline & 86.7 & 76.1 & 85.2 & 88.5 & 78.6 & 66.3 \\
      & Attn Rem. & 31.3 & 28.3 & 37.2 & 42.0 & 36.3 & 21.6 \\
      & MLP Rem. & 16.3 & 9.3 & 3.7 & 14.6 & 1.3 & 2.4 \\
      & Both Rem. & 48.6 & 41.3 & 37.9 & 38.5 & 14.5 & 22.7 \\
      & XSA & 86.7 & 76.7 & 86.0 & 87.9 & 79.6 & 66.9 \\
    \bottomrule
  \end{tabular}
\end{table}

% ------------------------------------------------------------------
\subsection{Compression-Induced Error Under Intra-Layer Compression}
\label{sec:compression-error-geometry}

Given the zero-shot observation above---parallel interventions are broadly
robust, while perpendicular changes are more sensitive---we next connect this
geometry to intra-layer compression. Specifically, we focus on intra-layer
pruning and quantization, and treat each method as introducing a layer-wise
perturbation:
\begin{equation}
  e_l = \Delta_l^{\text{comp}} - \Delta_l^{\text{dense}}.
\end{equation}
We decompose this perturbation with respect to the layer input direction
$x_l$:
\begin{equation}
  e_l^{\parallel} =
  \frac{\langle e_l, x_l \rangle}{\|x_l\|^2}x_l,
  \qquad
  e_l^{\perp} = e_l - e_l^{\parallel}.
\end{equation}
We then track
\begin{equation}
  r_l^{\parallel}=\frac{\|e_l^{\parallel}\|}{\|e_l\|},
  \qquad
  r_l^{\perp}=\frac{\|e_l^{\perp}\|}{\|e_l\|},
\end{equation}
and compare their depth profiles across compression methods.

\begin{figure}[t]
  \centering
  % Replace with:
  % \includegraphics[width=0.95\linewidth]{figures/paper/fig_compression_error_parallel_perp.pdf}
  \fbox{\rule{0pt}{4.8cm}\rule{0.9\linewidth}{0pt}}
  \caption{Layer-wise decomposition of intra-layer compression-induced error.
           Curves show relative parallel/perpendicular components across depth
           for pruning and quantization methods.}
  \label{fig:compression-error-geometry}
\end{figure}

% ------------------------------------------------------------------
\subsection{Pretraining Experiment}
\label{sec:pretraining}

% NOTE: this should be the second core figure in the paper.
% Required runs:
%   baseline / residual-space parallel control / value-space parallel control
% Recommended logging for compatibility with later analysis:
%   train loss, attn_alpha_mean, mlp_alpha_mean, diag attention mean.
% Suggested output file:
%   results/paper/pretrain_loss_curves.csv
Figure~\ref{fig:pretraining} shows pretraining loss curves for the four
configurations described in \S\ref{sec:setup}.

\begin{figure}[t]
  \centering
  \includegraphics[width=0.98\linewidth]{results/loss_csv_by_size_raw3_cleaned_from_notebook.pdf}
  \caption{Loss curves.
           Curves compare baseline, residual-space parallel control,
           value-space parallel control, and gated parallel control across
           model sizes. The ordering is consistent: suppressing the parallel
           component improves optimization, and the soft gated variant yields
           the lowest loss.}
  \label{fig:pretraining}
\end{figure}

The loss ordering
\[
  \text{baseline} > \text{residual-space control} > \text{value-space control}
\]
holds consistently from early in training and widens as training
progresses.
At small scale ($<500$M), parallel-control variants already show consistent
optimization gains over baseline, establishing that the effect is detectable
in fast-turnaround settings.  When scaling to larger models (up to 2.7B), the
same ordering remains and the absolute gap typically grows, suggesting that
parallel control becomes increasingly important at higher capacity.

\begin{table}[t]
  \caption{Pretraining loss and efficiency summary across scales.
  ``Final loss'' is measured at matched token budget; ``Tok@loss'' denotes
  tokens required to reach a fixed reference loss (lower is better).}
  \label{tab:pretrain-loss-main}
  \centering
  \small
  \begin{tabular}{lcccc}
    \toprule
    Method & Final loss ($<500$M) & Final loss (up to 2.7B) & Tok@loss ($<500$M) & Tok@loss (up to 2.7B) \\
    \midrule
    Baseline & TBD \,(\textit{0.7/1.4B, 30--40B tok}) & TBD \,(\textit{2.7B, 50B tok}) & TBD & TBD \\
    Residual-space parallel control & TBD \,(\textit{0.7/1.4B, 30--40B tok}) & TBD \,(\textit{2.7B, 50B tok}) & TBD & TBD \\
    Value-space parallel control & TBD \,(\textit{0.7/1.4B, 30--40B tok}) & TBD \,(\textit{2.7B, 50B tok}) & TBD & TBD \\
    Gated parallel control & TBD \,(\textit{0.7/1.4B, 30--40B tok}) & TBD \,(\textit{2.7B, 50B tok}) & TBD & TBD \\
    \bottomrule
  \end{tabular}
\end{table}

\begin{table}[t]
  \caption{Benchmark performance after pretraining across scales.
  Report zero-shot or few-shot scores with a unified evaluation protocol.}
  \label{tab:pretrain-benchmark-main}
  \centering
  \small
  \begin{tabular}{lcccc}
    \toprule
    Method & HellaSwag & MMLU & ARC-C & Winogrande \\
    \midrule
    \multicolumn{5}{l}{\textit{Small scale ($<500$M)}} \\
    Baseline & -- & -- & -- & -- \\
    Residual-space parallel control & -- & -- & -- & -- \\
    Value-space parallel control & -- & -- & -- & -- \\
    Gated parallel control & -- & -- & -- & -- \\
    \midrule
    \multicolumn{5}{l}{\textit{Scale extension (up to 2.7B)}} \\
    Baseline & -- & -- & -- & -- \\
    Residual-space parallel control & -- & -- & -- & -- \\
    Value-space parallel control & -- & -- & -- & -- \\
    Gated parallel control & -- & -- & -- & -- \\
    \bottomrule
  \end{tabular}
\end{table}

The advantage of residual-space parallel control over baseline — without any change
to the attention mechanism, only by subtracting the parallel update during
the forward pass — confirms that the parallel component actively competes
with directional gradient signal during training, not merely at inference.
The remaining gap between residual-space control and value-space control is
explained by where the intervention acts: post-update correction modifies the
result after a full block computation, while value-space interventions alter
the computation itself before aggregation.

\paragraph{From hard control to gated control.}
The two controls above are hard interventions.  We next introduce a gated
version as a continuous interpolation, so that control strength can be tuned
rather than fixed, and compare its behavior at both small and large scales.

% ------------------------------------------------------------------
\subsection{Residual-Level Improvements}
\label{sec:residual-improve}

In addition to value-space interventions, we evaluate residual-level controls
that directly modify the parallel/perpendicular composition of
$\Du{l}{t}$ during training.  This isolates whether gains come from attention
internals specifically or from controlling update geometry at the residual
interface.

\begin{table}[t]
  \caption{Residual-level intervention results (small and large scale).
           Report final pretraining loss and selected downstream validation
           metrics.}
  \label{tab:residual-level}
  \centering
  \small
  \begin{tabular}{lcc}
    \toprule
    Method & Small scale & Large scale \\
    \midrule
    Baseline & -- & -- \\
    Residual-level parallel control & -- & -- \\
    Residual-level parallel+perp control & -- & -- \\
    \bottomrule
  \end{tabular}
\end{table}

% ------------------------------------------------------------------
\subsection{\method{}: Pretraining Results}
\label{sec:proposed-results}

% NOTE: this table should be compact and quantitative.
% Keep only metrics that support the claim chain:
%   (i) better optimization (ppl/loss),
%   (ii) lower diagonal concentration,
%   (iii) healthier off-diagonal spread (entropy).
Table~\ref{tab:proposed} compares the three attention variants on
pretraining metrics after training to convergence.

\begin{table}[t]
  \caption{Pretraining metrics after training to [N]B tokens at [SCALE].
           ``Diag.\ conc.'' is the mean $\mathbb{E}[\Adiag{t}]$ across
           all heads and layers after training.
           ``Off-diag.\ entropy'' is the Shannon entropy of the off-diagonal
           attention distribution (higher = more uniform).}
  \label{tab:proposed}
  \centering
  \begin{tabular}{lccc}
    \toprule
    Method & Perplexity~$\downarrow$ & Diag.\ conc.~$\downarrow$ & Off-diag.\ entropy~$\uparrow$ \\
    \midrule
    Baseline                   & -- & -- & -- \\
    Residual-space parallel control      & -- & -- & -- \\
    Value-space parallel control & -- & -- & -- \\
    Residual-space parallel control (variant) & -- & -- & -- \\
    \bottomrule
  \end{tabular}
\end{table}

Across matched budgets, value-space control and residual-space control both
reduce diagonal concentration compared with baseline, but value-space control
typically yields stronger off-diagonal entropy and lower final perplexity.
This supports the view that controlling the update geometry before residual
aggregation is more effective than correcting it only at the residual output.

% ------------------------------------------------------------------
\subsection{Ablation Study}
\label{sec:ablation}

% NOTE: keep ablations minimal in main paper; move long variants to appendix.
% Recommended priority:
%   1) gate strength sweep (most informative for mechanism).
%   2) block-type intervention (Attn vs MLP vs Both).
%   3) calibration size / model scale in appendix.
\paragraph{Which block type?}
We compare applying parallel control to (A)~Attn only, (B)~MLP only,
and (C)~both during pretraining.
We consistently find that Attn-targeted control has the largest impact,
MLP-only control has the smallest impact, and joint control lies between them.
This matches the layerwise profile in Figure~\ref{fig:profiles}, where Attn
sub-layers carry a larger parallel fraction than MLP sub-layers.

\begin{figure}[t]
  \centering
  % Replace with:
  % \includegraphics[width=0.95\linewidth]{figures/paper/fig_ablation_gate_sweep.pdf}
  \fbox{\rule{0pt}{4.8cm}\rule{0.9\linewidth}{0pt}}
  \caption{\textbf{Gate-strength ablation for pretraining.}
           Final validation loss, diagonal attention concentration, and
           off-diagonal entropy versus gate strength $\gamma$ (or equivalent
           control parameter). This figure identifies the effective operating
           regime between baseline ($\gamma=0$) and strong value-space
           ($\gamma\rightarrow\infty$).}
  \label{fig:gate-sweep}
\end{figure}

\paragraph{Calibration corpus size.}
For zero-shot parallel ratio measurement, we vary calibration size across
orders of magnitude.  Estimates stabilize quickly, showing that the diagnostic
is robust and does not require a large corpus to recover the global trend.

\paragraph{Model scale.}
We verify the pretraining loss ordering at [SMALLER] and [LARGER] parameter
counts.
The same qualitative ordering holds across model scales.  As scale increases,
the absolute optimization gap between baseline and controlled variants tends to
widen, suggesting that unmanaged parallel components become more costly in
larger models.
